\documentclass[10pt,a4paper]{amsart}
\usepackage[margin=19mm]{geometry}
\usepackage{amsmath,amssymb,mathtools}
\usepackage[scr=rsfs,cal=euler]{mathalfa}
\usepackage{xfrac}
\usepackage[unicode,hidelinks,bookmarks=false]{hyperref}
\newcommand{\ie}{i.e.\ }
\newcommand{\cf}{cf.\ }
\newcommand{\resp}{respectively\ }
\newcommand{\Subsection}{\subsection}
\providecommand{\nobreakdash}{\nobreak\hbox{-}}
\setlength{\emergencystretch}{2em}
\multlinegap0pt
\usepackage{bm}
\usepackage{bbm}
\usepackage{dsfont}
\usepackage{xspace}
\usepackage{cancel}
\usepackage{mathtools}
\usepackage{amsthm}
\makeatletter
\@ifundefined{theorem}{\newtheorem{theorem}{Theorem}}{}
\@ifundefined{lemma}{\newtheorem{lemma}[theorem]{Lemma}}{}
\@ifundefined{proposition}{\newtheorem{proposition}[theorem]{Proposition}}{}
\makeatother
\newcommand{\E}{\mathbb{E}}
\newcommand{\Hqprime}{H^{(q) \prime}_{K,S}}
\newcommand{\Hq}{H^{(q)}_{K,S}}
\newcommand{\delx}{\partial_x}
\newcommand{\diff}{\mathrm{d}}
\newcommand{\ettaz}{\frac{S}{q} + \frac{K}{q+K}\frac{\mu - S}{q}}
\newcommand{\e}{\mathrm{e}}
\newcommand{\ind}{\mathds{1}}
\renewcommand{\geq}{\geqslant}
\title[A new perspective on the L{\o}kka-Zervos dichotomy for absol]{A new perspective on the L{\o}kka-Zervos dichotomy for absolutely continuous dividend strategies}
\author{Tommy Mastromonaco, Nacer Fendri, Jean-Fran\c{c}ois Renaud, Clarence Simard}
\date{Source: arXiv:2604.13302v4 (CC BY 4.0)}
\begin{document}
\maketitle

\noindent\emph{Source and licence.} A new perspective on the L{\o}kka-Zervos dichotomy for absolutely continuous dividend strategies. Version arXiv:2604.13302v4 (CC BY 4.0), distributed under the Creative Commons Attribution 4.0 International licence, as recorded in the arXiv licence metadata for this exact version and shown on the abstract page https://arxiv.org/abs/2604.13302v4. Full rights notice: licenses/arxiv-2604.13302-CC-BY-4.0.txt.\par\smallskip

\noindent\textit{Editorial scope.} Counted unit: the Lemma whose source label is \texttt{lem-delxJd} in the article (source0.tex lines 406--413, source label \texttt{lem-delxJd}), delivered together with the article's own complete original proof of it (source0.tex lines 414--434, one \texttt{proof} environment of 21 lines). No mathematics is written, repaired or re-derived: statement and proof are the article's own text line for line. Required context retained uncounted: thm:sol-Jd (lines 208--230); cond-params-bd (lines 210--212); thm:sol-Jc (lines 243--253); Jd-compact (lines 345--347); Jc-compact (lines 349--351); prop:seuils-croiss (lines 396--398). Every definition, display and label of those lines is retained unchanged. Omitted from this package and not redistributed: 1--207, 213--242, 254--344, 348--348, 352--395, 399--405, 435--688. Nothing inside a retained excerpt is omitted or altered: every retained body is delivered exactly as the article's lines are written, so no omission receipt of any kind accompanies this package. Citations: the retained excerpts carry no citation command at all, so no bibliography entry of the article is transcribed and no reference is redistributed. Reference adaptation: none was needed and none was made; every reference inside the delivered unit resolves inside it, so no unresolved reference or citation is produced. Printed slips kept literal: no printed word, symbol, hypothesis or number of the counted statement or of its proof was altered. Layout and class: the article's own class is \texttt{amsart} with packages including a4wide, babel, hyperref, amsfonts, amsmath, amssymb, amsthm, amsaddr, latexsym, layout, dsfont, xcolor, graphicx, mathtools; the wrapper is the lane's pinned \texttt{amsart} editorial wrapper at 10pt on A4, which restates the theorem-like environments the retained text needs and adds the article's own macro definitions that the retained text needs (\texttt{\textbackslash E}, \texttt{\textbackslash Hq}, \texttt{\textbackslash Hqprime}, \texttt{\textbackslash delx}, \texttt{\textbackslash diff}, \texttt{\textbackslash e}, \texttt{\textbackslash ettaz}, \texttt{\textbackslash geq}, \texttt{\textbackslash ind}), with extra wrapper packages (bm, bbm, dsfont, xspace, cancel, mathtools). The delivered numbering is the unit's own. Assets: no retained span contains a figure environment, a graphics-inclusion command, a PDF-page inclusion, a table or a TikZ picture; no image file is copied into this package, the package declares no asset dependency and no image file is redistributed with it. Licence: version arXiv:2604.13302v4 of this article is distributed under the Creative Commons Attribution 4.0 International licence, as recorded in the arXiv licence metadata for this exact version and shown on the abstract page https://arxiv.org/abs/2604.13302v4. A full rights notice is delivered at licenses/arxiv-2604.13302-CC-BY-4.0.txt. Characters: every retained excerpt is pure ASCII, so the delivered engine, XeLaTeX with its default Latin Modern text font, reports no missing character.
\begin{theorem}\label{thm:sol-Jd}
    If
    \begin{equation}\label{cond-params-bd}
        \frac{q}{q+K}\frac{\Hq(0)}{\Hqprime(0)} + \ettaz + P > 0,
    \end{equation}
    then there exists a unique $b_d>0$ such that
    \begin{equation}\label{Vd-proba}
        V_d(x) = \E_x \left[ \int_0^{\tau_d} \mathrm{e}^{-qt} (KX^{b_d}_t +S) \ind_{X^{b_d}_t \geq b_d} \diff t - P\e^{-q\tau_d}\right]        
    \end{equation}
    with $\tau_d = \inf\{t\geq 0\mid X^{b_d}_t < 0\}$, where $X^{b_d}$ is the solution to
    \begin{equation}
        \diff X^{b_d}_t = \left(\mu - (K X^{b_d}_t + S) \ind_{X^{b_d}_t \geq b_d} \right) \diff t + \sigma \diff B_t .
    \end{equation}
    
    Otherwise, if the condition in Equation~\eqref{cond-params-bd} fails, then we have
    \begin{equation}\label{Vd-bd0-proba}
        V_d(x) = \E_x \left[ \int_0^{\tau_d} \mathrm{e}^{-qt} (KX'_t +S) \diff t - P\e^{-q\tau_d}\right]
    \end{equation}
    with $\tau_d = \inf\{t\geq 0\mid X'_t < 0\}$, where $X'$ is the solution to
    \begin{equation}
        \diff X'_t = \left(\mu - (K X'_t + S) \right) \diff t + \sigma \diff B_t .
    \end{equation}
\end{theorem}

\begin{theorem}\label{thm:sol-Jc}
    There exists a unique $b_c>0$ such that
    \begin{equation}\label{Vc-proba}
        V_c(x) = \E_x \left[\int_0^\infty \e^{-qt}\left((K Y^{b_c}_t + S)\ind_{Y^{b_c}_t\geq b_c} \diff t - \beta\diff G^{b_c}_t\right)\right],        
    \end{equation}
    where $(Y^{b_c},G^{b_c})$ is the solution to the Skorokhod problem
    \begin{align*}
        \diff Y^{b_c}_t &= \left(\mu - (K Y^{b_c}_t + S)\ind_{Y^{b_c}_t\geq b_c}\right)\diff t + \sigma \diff B_t + \diff G^{b_c}_t ,\\
            G^{b_c}_t &= \int_0^t \ind_{Y^{b_c}_s = 0} \diff G^{b_c}_s .
    \end{align*}
\end{theorem}

\begin{equation}\label{Jd-compact}
J_d(x;b) = R(x;b) - P\Phi(x;b)
\end{equation}

\begin{equation}\label{Jc-compact}
    J_c(x;b) = R(x;b) + J_c(0;b)\Phi(x;b) .
\end{equation}

\begin{proposition}\label{prop:seuils-croiss}
    Suppose  $P^\prime$ satisfies equation \eqref{cond-params-bd} with $P<P^\prime$, then $b_d(P)<b_d(P^\prime)$. Suppose $1< \beta < \beta^\prime$, then $b_c(\beta)<b_c(\beta^\prime)$.
\end{proposition}

\begin{lemma}\label{lem-delxJd}
    We have that:
    \begin{enumerate}
        \item $\delx J_d(0;b_d) \geq \delx J_d(0;b)$, for all $b\geq 0$;
        \item for $b\geq 0$, $\delx J_d(0;b) =\beta$ if and only if $J_c(0;b) = -P$;
        \item if $J_c(0;b_c) < -P$ (resp.\ $>$ or $=$), then $b_d < b_c$ (resp.\ $>$ or $=$).
    \end{enumerate}
\end{lemma}

\begin{proof}
Since $V_d = J_d(\cdot;b_d)$ and since, for any $b\geq 0$, $J_d(\cdot;b)$ is the performance function of a strategy $\pi^b \in \Pi_d$, then we have $J_d(x;b_d) \geq J_d(x;b)$ for all $x \geq 0$ and any $b\geq 0$. As $J_d(0;b)=-P$ for any $b \geq 0$, then we have $\delx J_d(0;b_d) \geq \delx J_d(0;b)$. This proves the first assertion.
    
Now, using the unified analytical representations of $J_d$ and $J_c$ (see~\eqref{Jd-compact} and~\eqref{Jc-compact}), and using the fact that $\delx J_c(0;b)=\beta$ for any $b \geq 0$, we have
\begin{equation}
\delx J_d(0;b) = \delx R(0;b) - P \delx\Phi(0;b)
\end{equation}
and
\begin{equation}
\delx J_c(0;b) = \delx R(0;b) + J_c(0;b) \delx\Phi(0;b) = \beta ,
\end{equation}
from which the second result follows.
    
For the last assertion, let us prove first the case corresponding to an equality, i.e., if $J_c(0;b_c) = -P$, then $b_c=b_d$. Suppose that $J_c(0;b_c) = -P$. Then, by the second assertion, we have $\delx J_d(0;b_c) =\beta$. Note from Theorem~\ref{thm:sol-Jd} and Theorem~\ref{thm:sol-Jc} that $V^\prime_d(b_d) = 1 = V^\prime_c(b_c)$, so that $b_d$ and $b_c$ uniquely satisfy $\delx J_d(b_d;b_d) = 1 = \delx J_c(b_c;b_c)$. Thus,
\begin{equation}
\delx R(b_d;b_d) - P\delx\Phi(b_d;b_d) = \delx R(b_c;b_c) - P\delx\Phi(b_c;b_c),  
\end{equation}
where we used the assumption that $J_c(0;b_c) = -P$. Clearly, $b_d = b_c$ solves this equation.
    
Suppose now that $J_c(0;b_c) < -P$ and set $P^\prime := -J_c(0;b_c) > P$. Then, by the result shown in the previous paragraph, we have $b_d(P^\prime)=b_c > 0$. Furthermore, we have $b_d(P^\prime) > b_d(P)$ by Proposition~\ref{prop:seuils-croiss}. Consequently, $b_c > b_d(P)$, that is, $b_d<b_c$. The same arguments hold with reverse inequalities in the case $J_c(0;b_c) > -P$.
\end{proof}

\end{document}
